\section{Discussion}
\label{sec:discussion}

\paragraph{What a static tie does and does not establish.}
Static preference accuracy remains useful: it measures correctness on a fixed,
independent evaluation distribution. However, it is a many-to-one summary of
behavior. The Attribute Phi--LLaVA tie and the complete matched-pair census show
that a shared measured score need not determine how a judge responds when the
visual fact supporting a preference changes. RewardLens is therefore a
complement to, rather than a replacement for, static evaluation: it asks for a
second measurement axis when behavioral response to visual changes matters.

\paragraph{Adaptation is distinct from generic stability.}
The two intervention branches prevent a simple preference for either changing
or retaining a prediction. On a relevant edit, retention by a base-correct
judge is an error because the gold preference changes; on an irrelevant edit,
retention is the desired response. Qwen's Count behavior makes this distinction
concrete: its high II does not offset its low RA. Reporting both quantities
therefore distinguishes selective adaptation from generic prediction stability.

\paragraph{Scope.}
RewardLens measures observable behavior under four controlled visual factors;
it does not identify a model's internal reasoning process or establish causal
claims about deployed-system outcomes. The results should consequently be
interpreted as evidence that static scores can leave factor-specific
intervention behavior unresolved in this evaluation setting, not as a universal
ranking of multimodal judges.
